\chapter*{Introduction générale}
\addcontentsline{toc}{chapter}{Introduction générale}

Dans le contexte actuel de transformation numérique, la gestion de l'information et de la documentation constitue un enjeu stratégique pour toutes les organisations. La capacité à produire des documents fiables, cohérents et personnalisés, tout en garantissant la sécurité des accès, est devenue une nécessité opérationnelle. Pourtant, de nombreux processus administratifs reposent encore sur des manipulations manuelles de fichiers bureautiques, ce qui expose l'organisation à des erreurs de saisie, à des problèmes de versioning et à une hétérogénéité visuelle.

La problématique centrale de ce projet de fin d'études réside dans la conception et la réalisation d'une solution capable d'automatiser le cycle de vie documentaire. Il s'agit de passer d'une gestion statique, basée sur des fichiers isolés, à une gestion dynamique, où le document est le résultat d'une association contrôlée entre des modèles éditables, des familles de documents, des chartes graphiques et des données métier extraites directement du système d'information.

L'objectif de ce projet est donc de fournir une plateforme web complète et sécurisée permettant aux administrateurs de définir des structures documentaires flexibles et aux utilisateurs finaux de générer des documents personnalisés en quelques clics. La solution s'appuie sur une architecture moderne utilisant ASP.NET Core pour le backend, Angular pour le frontend et SQL Server pour la persistance des données.

Le présent rapport décrit la démarche de réalisation de cette plateforme, organisée en six chapitres :
\begin{itemize}
    \item le premier chapitre présente le cadre général du projet, l'organisme d'accueil et l'environnement technique ;
    \item le deuxième chapitre détaille l'analyse des besoins et la spécification fonctionnelle ;
    \item les chapitres trois à six décrivent la réalisation technique à travers les quatre sprints du projet, couvrant l'authentification, la gouvernance d'accès, l'édition documentaire, puis l'administration métier avec l'intégration finale.
\end{itemize}

L'ensemble de ce travail vise à démontrer comment une approche logicielle structurée peut répondre concrètement aux défis de la gestion documentaire moderne.

